% ============================================================
\chapter{Lyapunov Exponents}
\label{ch:exponents}
% ============================================================

\section{Introduction}

How does one measure chaos? Sensitivity to initial conditions is a qualitative notion: two nearby orbits diverge, but how fast? Alexander Lyapunov, in his 1892 thesis, had already introduced characteristic exponents to study the stability of differential equations. But it was in the 1960s--1970s, with the work of Valery Oseledets (the multiplicative ergodic theorem, 1968), that these exponents found their definitive formulation for general dynamical systems. The maximal Lyapunov exponent quantifies the rate of exponential divergence of nearby orbits: if it is positive, the system is chaotic. The full spectrum of exponents reveals the geometry of the attractor, distinguishing directions of expansion, contraction, and neutrality.

\begin{intuition}
Imagine two particles very close together in a flow. If they separate
exponentially fast, the Lyapunov exponent is positive and the system is
chaotic. If they converge, the exponent is negative and the system is stable.
The exponent measures precisely this exponential rate.
\end{intuition}

\section{Definition for Discrete Systems}

\begin{definition}[Lyapunov exponent for a one-dimensional map]
\label{def:lyap-1d}
\index{Lyapunov exponent!one-dimensional}
Let $f : \R \to \R$ be $\mathcal{C}^1$ and $x_0$ an initial point. The
\textbf{Lyapunov exponent} at $x_0$ is
\[
    \lambda(x_0) = \lim_{n \to \infty} \frac{1}{n}
    \sum_{k=0}^{n-1} \ln \abs{f'(x_k)},
\]
where $x_k = f^k(x_0)$, whenever this limit exists.
\end{definition}

\begin{remark}
By the chain rule, $\abs{(f^n)'(x_0)} = \prod_{k=0}^{n-1}\abs{f'(x_k)}$,
so $\lambda(x_0) = \lim_{n \to \infty} \frac{1}{n}\ln\abs{(f^n)'(x_0)}$.
The number $e^\lambda$ is the average stretching factor per iteration.
\end{remark}

\begin{example}[Logistic map]
For $f_4(x) = 4x(1-x)$, the invariant measure is
$\mu(\mathrm{d}x) = \frac{1}{\pi\sqrt{x(1-x)}}\mathrm{d}x$. The Lyapunov
exponent is
\[
    \lambda = \int_0^1 \ln\abs{4 - 8x}\,\frac{\mathrm{d}x}{\pi\sqrt{x(1-x)}}
    = \ln 2.
\]
\end{example}

\section{Higher-Dimensional Lyapunov Exponents}

\begin{definition}[Lyapunov matrix]
\index{Lyapunov matrix}
Consider $\dot{x} = f(x)$ in $\R^n$ with fundamental matrix $\Phi(t, x_0)$
of the variational equation $\dot{Y} = Df(\varphi(t, x_0))Y$. The
\textbf{Lyapunov matrix} is
\[
    \Lambda(x_0) = \lim_{t \to \infty}
    \left[\Phi(t, x_0)^\top \Phi(t, x_0)\right]^{1/(2t)}.
\]
\end{definition}

\begin{definition}[Lyapunov spectrum]
\index{Lyapunov spectrum}
The \textbf{Lyapunov exponents} of the system at $x_0$ are the logarithms of
the eigenvalues of $\Lambda(x_0)$, denoted
\[
    \lambda_1 \geq \lambda_2 \geq \cdots \geq \lambda_n.
\]
They measure the rates of expansion in the $n$ principal directions.
\end{definition}

\begin{keyformulas}[title=\textbf{Geometric Interpretation of the Exponents}]
\begin{itemize}
    \item $\lambda_i > 0$: exponential stretching in the $i$-th direction.
    \item $\lambda_i = 0$: neither contraction nor expansion (flow direction).
    \item $\lambda_i < 0$: exponential contraction.
    \item $\lambda_1 > 0$: the system is \textbf{chaotic}.
    \item $\sum_{i=1}^n \lambda_i < 0$: the system is \textbf{dissipative}.
\end{itemize}
\end{keyformulas}

\section{Oseledets Multiplicative Ergodic Theorem}

\begin{theorem}[Oseledets, 1968]
\label{thm:oseledets}
\index{Oseledets theorem}
Let $f : M \to M$ be a diffeomorphism preserving an ergodic probability
measure $\mu$, and let $A(x) = Df(x)$. Suppose
$\int \ln^+ \norm{A(x)}\,\mathrm{d}\mu(x) < \infty$. Then, for $\mu$-almost
every $x$, there exist real numbers
$\lambda_1 > \lambda_2 > \cdots > \lambda_s$ and a decomposition
\[
    \R^n = E_1(x) \oplus E_2(x) \oplus \cdots \oplus E_s(x)
\]
such that for every $v \in E_i(x) \setminus \{0\}$:
\[
    \lim_{n \to \pm\infty} \frac{1}{n}\ln\norm{Df^n(x) \cdot v} = \lambda_i.
\]
Moreover, $Df(x) \cdot E_i(x) = E_i(f(x))$ and $\dim E_i(x)$ is constant
$\mu$-almost everywhere.
\end{theorem}

\begin{remark}
Oseledets' theorem generalises Lyapunov exponents to any linear cocycle over
an ergodic system. The integers $m_i = \dim E_i(x)$ are the
\textbf{multiplicities} of the exponents.
\end{remark}

\section{Fundamental Properties}

\begin{proposition}[Exponent along the flow]
For an autonomous continuous system $\dot{x} = f(x)$, the flow direction
$f(x)$ is always associated with a zero Lyapunov exponent. For a chaotic
attractor in dimension $n$, the spectrum satisfies
$\lambda_1 > 0 = \lambda_2 > \lambda_3 \geq \cdots \geq \lambda_n$.
\end{proposition}

\begin{proposition}[Liouville formula]
The sum of Lyapunov exponents satisfies
\[
    \sum_{i=1}^{n} \lambda_i = \lim_{T \to \infty} \frac{1}{T}
    \int_0^T \tr\bigl(Df(\varphi(t, x_0))\bigr)\,\mathrm{d}t.
\]
For the Lorenz system ($\sigma=10$, $\rho=28$, $\beta=8/3$):
$\lambda_1 + \lambda_2 + \lambda_3 = -(\sigma + 1 + \beta) \approx -13.67$.
\end{proposition}

\begin{theorem}[Pesin's formula]
\index{Pesin formula}
Let $f$ be a $\mathcal{C}^2$ diffeomorphism of a compact manifold preserving
an ergodic measure $\mu$ absolutely continuous with respect to Lebesgue. Then
\[
    h_\mu(f) = \sum_{\lambda_i > 0} m_i\,\lambda_i,
\]
where the sum is over positive exponents counted with multiplicity.
\end{theorem}

\begin{attention}[title=\textbf{Exponents and dimension}]
The Kaplan-Yorke formula relates Lyapunov exponents to the attractor
dimension: $D_{KY} = j + \frac{\sum_{i=1}^{j}\lambda_i}{\abs{\lambda_{j+1}}}$,
where $j$ is the largest integer with $\sum_{i=1}^{j}\lambda_i \geq 0$. This
is conjectured to equal the information dimension.
\end{attention}

\section{Numerical Computation}

\begin{definition}[Benettin's algorithm]
\index{Benettin algorithm}
The standard algorithm for computing the first $p$ Lyapunov exponents:
\begin{enumerate}
    \item Initialise an orthonormal basis $\{v_1, \ldots, v_p\}$.
    \item Simultaneously integrate the orbit and $p$ tangent vectors over an
    interval $\tau$.
    \item Apply Gram-Schmidt to re-orthogonalise. Record the norms before
    orthogonalisation.
    \item Repeat steps 2--3 for $N$ iterations.
    \item Estimate $\lambda_i \approx \frac{1}{N\tau}\sum_{k=1}^{N}\ln\norm{v_i^{(k)}}$.
\end{enumerate}
\end{definition}

\begin{codeblock}[title=\textbf{Lyapunov exponents of the Lorenz system}]
\begin{minted}{python}
import numpy as np
from scipy.integrate import solve_ivp

def lorenz(t, state, sigma=10, rho=28, beta=8/3):
    x, y, z = state[:3]
    dxdt = [sigma*(y - x), rho*x - y - x*z, x*y - beta*z]
    J = np.array([[-sigma, sigma, 0],
                  [rho - z, -1, -x],
                  [y, x, -beta]])
    Y = state[3:].reshape(3, 3)
    dYdt = J @ Y
    return dxdt + dYdt.flatten().tolist()

n = 3
x0 = [1.0, 1.0, 1.0] + list(np.eye(n).flatten())
dt, N_steps = 0.01, 10000
lyap = np.zeros(n)

for i in range(N_steps):
    sol = solve_ivp(lorenz, [0, dt], x0, max_step=dt/2)
    x0[:3] = sol.y[:3, -1]
    Y = sol.y[3:, -1].reshape(n, n)
    Q, R = np.linalg.qr(Y)
    lyap += np.log(np.abs(np.diag(R)))
    x0[3:] = Q.flatten()

lyap /= (N_steps * dt)
print(f"Lyapunov exponents: {lyap}")
# Typical result: [0.91, 0.00, -14.57]
\end{minted}
\end{codeblock}

\section{Applications to Chaos Detection}

\begin{proposition}[Chaos criterion]
A deterministic system is chaotic if and only if its largest Lyapunov exponent
is strictly positive: $\lambda_1 > 0$.
\end{proposition}

\begin{example}[Classification by Lyapunov exponents]
For a three-dimensional system:
\begin{itemize}
    \item $(-, -, -)$: stable fixed point.
    \item $(0, -, -)$: stable limit cycle.
    \item $(0, 0, -)$: quasi-periodic torus $T^2$.
    \item $(+, 0, -)$: strange attractor (chaos).
\end{itemize}
\end{example}

\begin{lemma}[Predictability time]
The Lyapunov time $\tau_L = 1/\lambda_1$ defines the characteristic time
scale beyond which prediction becomes impossible. After time $t$, an initial
uncertainty $\delta_0$ grows as $\delta(t) \sim \delta_0\,e^{\lambda_1 t}$.
\end{lemma}

\begin{example}[Weather prediction]
For the atmosphere, $\lambda_1 \approx 1/(2\text{ days})$. An initial error
of $10^{-6}$ reaches the synoptic scale ($\sim 1$) after roughly
$t = \frac{\ln(10^6)}{\lambda_1} \approx 28$ days, setting the theoretical
limit of weather prediction.
\end{example}

\section{Conditional Exponents and Synchronisation}

\begin{definition}[Conditional Lyapunov exponents]
\index{conditional Lyapunov exponents}
For coupled systems $\dot{x} = f(x)$ and $\dot{y} = g(x, y)$, the
\textbf{conditional Lyapunov exponents} of the slave system $y$ are defined
via $\dot{\delta y} = D_y g(x(t), y(t))\cdot\delta y$. They determine whether
synchronisation $y(t) \to x(t)$ is stable.
\end{definition}

\begin{theorem}[Chaotic synchronisation]
Two identical coupled chaotic systems synchronise if and only if all
transverse conditional Lyapunov exponents are strictly negative.
\end{theorem}

\section{Exercises}

\begin{exercise}[Tent map exponent]
Analytically compute the Lyapunov exponent of the tent map
$T(x) = 1 - \abs{2x - 1}$ on $[0,1]$ with respect to Lebesgue measure.
\end{exercise}

\begin{exercise}[Arnold's cat map spectrum]
For Arnold's cat map
$\begin{pmatrix} x_{n+1} \\ y_{n+1} \end{pmatrix}
= \begin{pmatrix} 2 & 1 \\ 1 & 1 \end{pmatrix}
\begin{pmatrix} x_n \\ y_n \end{pmatrix} \pmod{1}$,
compute both Lyapunov exponents and verify that their sum is zero
(conservative system).
\end{exercise}

\begin{exercise}[Kaplan-Yorke dimension]
For the Lorenz system with standard parameters, the exponents are
approximately $\lambda_1 \approx 0.91$, $\lambda_2 = 0$,
$\lambda_3 \approx -14.57$. Compute the Kaplan-Yorke dimension.
\end{exercise}

\begin{exercise}[Numerical convergence]
Implement Benettin's algorithm for the H\'enon map
$H(x,y) = (1 - 1.4x^2 + y,\; 0.3x)$ and study the convergence of the
exponents as a function of the number of iterations.
\end{exercise}

\begin{exercise}[Pesin's formula]
Numerically verify Pesin's formula for Arnold's cat map by independently
computing the metric entropy and the sum of positive exponents.
\end{exercise}
